% _plantilla.tex -- plantilla OBLIGATORIA de capítulo (R7). NO se incluye en qfdesk.tex.
%
% Uso: copiar a capitulos/capNN_<tema>.tex, sustituir los textos entre «<…>» y borrar estos comentarios.
%  * Un capítulo = un archivo. Empieza con \chapter{<título>}\label{cap:<etiqueta>} (las etiquetas de los
%    capítulos 1-15 están fijadas en qfdesk.tex / README.md; los demás capítulos solo se referencian
%    con \ref{cap:...}, nunca con etiquetas internas de otro capítulo).
%  * Etiquetas internas con prefijo del capítulo: eq:capNN-…, sec:capNN-…, tab:capNN-….
%  * NO definir macros (\newcommand, \renewcommand, \def, \DeclareMathOperator): se usan las de
%    preambulo.tex (apéndice A). Un símbolo propio se escribe en LaTeX explícito, se declara en la sección
%    «Notación» y se añade a simbolos/capNN.tex.
%  * \todo solo en los capítulos marcador. Documentar lo que el código HACE, no la teoría del libro.
%  * Cuotas: >= 11 \codigo en cada capítulo 2-12 (>= 5 en 1 y 13), >= 2 \verif (capítulos 2-12, cada uno
%    con su script en verificacion/capNN_<tema>.py), >= 2 citas distintas (capítulos 2-12).
%
\chapter{<Título del capítulo>}\label{cap:<etiqueta>}

% -------------------------------------------------------------------------------------------------
\section{Propósito}\label{sec:capNN-proposito}
% 1. Qué calcula el código de este tema y para qué se usa en la aplicación (2-4 frases). Remitir a otros
%    capítulos con \ref{cap:curvas}, \ref{cap:swaps-ois}, etc.

<Qué calcula el código y para qué.>

% -------------------------------------------------------------------------------------------------
\section{Notación}\label{sec:capNN-notacion}
% 2. Solo los símbolos NUEVOS del capítulo (cada uno con su unidad); los símbolos comunes ($\DF{c}{T}$,
%    $\yf{t_1}{t_2}$, $\Nom$, $\K$, ...) remiten al apéndice \ref{ap:simbolos}. Añadir la fila de cada
%    símbolo nuevo a simbolos/capNN.tex.

Los símbolos comunes están en el apéndice~\ref{ap:simbolos}. En este capítulo, $w_k$ es el peso del vértice~$k$
(adimensional) y $\tau_i = \yf{t_{i-1}}{t_i}$ la fracción de año del periodo~$i$ (años).

% -------------------------------------------------------------------------------------------------
\section{Modelo y fórmulas}\label{sec:capNN-modelo}
% 3. Fórmulas numeradas (equation / align) con TODAS las convenciones explícitas: base de días,
%    composición, interpolación, signo, unidades, fecha de valuación. Cada fórmula corresponde a una
%    función concreta: pegar el \codigo junto a la fórmula. Dos ecuaciones de ejemplo (formato, no
%    contenido de ningún capítulo):

La tasa par de un swap OIS con fechas de pago $T_1,\dots,T_n$ y fecha de inicio $T_0$ es
\begin{equation}\label{eq:capNN-tasa-par}
  \Spar = \frac{\DF{c}{T_0} - \DF{c}{T_n}}{\Ann}, \qquad
  \Ann = \sum_{i=1}^{n} \yf{T_{i-1}}{T_i}\,\DF{c}{T_i},
\end{equation}
con $\yf{\cdot}{\cdot}$ en base Actual/360. Implementación:
\codigo{include/qf/instruments/ois_swap.hpp}{OisSwap}.

% Recuadro obligatorio cuando el código se aparta del libro (R3). Si no hay diferencia, decirlo en una
% frase («No hay diferencia con la referencia en este capítulo.»):
\begin{diferencia}
  <Qué hace el código distinto de la referencia: por ejemplo, el cupón del bono, la composición continua del
  periodo en curso o la base determinista. Citar la fuente con \citet{clave}.>
\end{diferencia}

% -------------------------------------------------------------------------------------------------
\section{Implementación}\label{sec:capNN-implementacion}
% 4. Mapa fórmula $\to$ función. Cada cálculo con al menos un \codigo{ruta}{símbolo}: la ruta es relativa
%    a la raíz del repositorio, el símbolo es una función, clase o constante (sin números de línea). Se
%    escribe con `_` normal (NO \_ ). referencias_codigo.py comprueba que la ruta existe y que el símbolo
%    está definido en el archivo.

La ecuación~\eqref{eq:capNN-tasa-par} se evalúa en C++ en
\codigo{include/qf/instruments/ois_swap.hpp}{OisSwap}; el binding Python está en
\codigo{src/python_bindings/qfpy.cpp}{par_rate} y el backend lo orquesta en
\codigo{desk/api/pricing/rates.py}{_irs_price}.

% -------------------------------------------------------------------------------------------------
\section{Ejemplo numérico verificable}\label{sec:capNN-ejemplo}
% 5. (>= 2 por capítulo 2-12.) Insumos concretos, fórmula aplicada y resultado con \verif{id}{valor}:
%    `id` = capNN-nombre, único en TODO el documento; `valor` en decimal con punto, sin separadores de
%    miles, con signo y exponente `e` opcionales y los dígitos suficientes. La tolerancia va en el texto y
%    es >= media unidad del último dígito impreso. Cada id se comprueba en exactamente un script
%    verificacion/capNN_<tema>.py que llama al código del proyecto (README.md, sección «Ejemplos»).

Con $c$ plana a $\K = 0.04$ (composición continua), $T_0 = 0$ y $T_n = 1$ año, $\DF{c}{T_n} = e^{-0.04} =
\verif{capNN-ejemplo-df}{0.9607894392}$; el resultado se reproduce con
\codigo{docs/tex/verificacion/capNN_<tema>.py}{<función>} con tolerancia absoluta $10^{-9}$.

% -------------------------------------------------------------------------------------------------
\section{Validación}\label{sec:capNN-validacion}
% 6. Pruebas del repositorio que cubren el capítulo (se citan por ruta con \codigo; símbolo = nombre de la
%    función de prueba) y, si aplica, los casos contra QuantLib del capítulo~\ref{cap:validacion}.

<Pruebas del repositorio y casos contra QuantLib que cubren este tema.>

% -------------------------------------------------------------------------------------------------
\section{Límites y supuestos}\label{sec:capNN-limites}
% 7. Los límites y supuestos que ya declara la Metodología de la aplicación (datos sintéticos,
%    aproximaciones, lo no modelado) y lo que el código no hace (R6). Nada se describe como hecho si el
%    código no lo hace. Discrepancias código <-> Metodología: NO se corrigen; van a HALLAZGOS.md.
\begin{limites}
  \begin{itemize}
    \item <Supuesto o límite 1.>
    \item <Supuesto o límite 2.>
  \end{itemize}
\end{limites}

% Solo capítulo 10 (coberturas): aviso de revisión pendiente por un contador IFRS 9.
% \begin{aviso}[Aviso IFRS 9]
%   <Texto del aviso.>
% \end{aviso}

% -------------------------------------------------------------------------------------------------
\section{Referencias}\label{sec:capNN-referencias}
% 8. Las citas viven dentro del texto (\citet{clave}, \citep{clave}); la bibliografía es única, al final
%    del documento (referencias.bib). Aquí, una frase que oriente al lector sobre las fuentes del capítulo.
%    Antes de añadir una obra nueva al .bib: confirmarla contra crossref.org u openlibrary.org
%    (libros y artículos con DOI o ISBN; normas e instituciones como @manual con URL oficial).

El tratamiento de las curvas sigue a \citet{hagan2006} y el modelo de tasas a \citet{brigo2006}.
